\begin{statementbox}

	\textbf{\textcolor{CStatement}{条件}}
	\begin{description}[style=nextline, leftmargin=2.8em, labelsep=0.5em, itemsep=0.35em]
		\item[\textbf{\textcolor{CStatement}{条件1（棱长a）：}}] 正方体棱长为$a$。
	\end{description}

	\textbf{\textcolor{CStatement}{结论}}
	\begin{description}[style=nextline, leftmargin=2.8em, labelsep=0.5em, itemsep=0.45em]
		\item[\textbf{\textcolor{CStatement}{结论一（内切球半径）：}}]
			\textbf{\textcolor{CStatement}{直观描述：}}
			内切球与六个面中心相切，半径等于棱长的一半。
			\[
				r_i = \frac{a}{2}.
			\]

		\item[\textbf{\textcolor{CStatement}{结论二（棱切球半径）：}}]
			\textbf{\textcolor{CStatement}{直观描述：}}
			棱切球与十二条棱中点相切，半径等于面对角线长的一半。
			\[
				r_m = \frac{\sqrt{2}}{2}a.
			\]

		\item[\textbf{\textcolor{CStatement}{结论三（外接球半径）：}}]
			\textbf{\textcolor{CStatement}{直观描述：}}
			外接球过八个顶点，半径等于体对角线长的一半。
			\[
				r_o = \frac{\sqrt{3}}{2}a.
			\]

		\item[\textbf{\textcolor{CStatement}{推广结论（体对角线关系）：}}]
			\textbf{\textcolor{CStatement}{直观描述：}}
			若正方体的体对角线长为$d$，则外接球半径也可表示为$d/2$。
			\[
				r_o = \frac{d}{2},\quad d = \sqrt{3}a.
			\]
	\end{description}

\end{statementbox}
